\subsection{Modelagem Cinemática}
\label{sec:modelagem_cinematica}

A modelagem cinemática descreve a relação geométrica entre os ângulos das juntas e a atitude (posição e orientação) do efetuador final. Utiliza-se a formulação clássica de Denavit--Hartenberg (DH), pois representa a transformação entre os referenciais das juntas consecutivas.

% --------------------------------------
% Explicação:
% A cinemática é dividida em:
%%%%% 1. Cinemática direta:
% dada a configuração articular (θ, d), obter posição/orientação do efetuador.
%%%%% 2. Cinemática inversa:
% dado o objetivo de posição/orientação do efetuador, calcular variáveis articulares.
% Começa pela cinemática direta, utilizando as matrizes homogêneas A_i da convenção DH.
% --------------------------------------
Cada junta do MR é associada a um sistema de coordenadas local, e a relação entre dois sistemas é descrita por:

\begin{equation}
A_i =
\begin{bmatrix}
\cos\theta_i & -\sin\theta_i \cos\alpha_i & \sin\theta_i \sin\alpha_i & a_i \cos\theta_i \\
\sin\theta_i & \cos\theta_i \cos\alpha_i & -\cos\theta_i \sin\alpha_i & a_i \sin\theta_i \\
0 & \sin\alpha_i & \cos\alpha_i & d_i \\
0 & 0 & 0 & 1
\end{bmatrix},
\end{equation}

onde, $\theta_i$ é o ângulo da junta revoluta ($i = 1,2,3,4$), $d_i$ é o deslocamento da junta prismática ($i = 5$); $a_i$ é o comprimento do elo $i$; e $\alpha_i$ é o ângulo de torção entre os eixos.
% ---------------------------------------------------------
% Comentário:
% A convenção DH padroniza os parâmetros geométricos.
% Em cada junta, definimos 4 parâmetros (θ, d, a, α).
% O produto sucessivo das matrizes A_i fornece a transformação total da base {0} até o efetuador {E}.
% ---------------------------------------------------------
A cinemática direta do MR é obtida pelas transformações:

\begin{equation}
    \label{eq:matrizes_homogeneas_sucessivas}
{}^0 \mathbf{T}_{E} = A_1 A_2 A_3 A_4 A_5,
\end{equation}

% Extração de orientação e posição a partir de \mathbf{T}^0_E:
% A orientação é a submatriz 3x3 superior-esquerda; 
% a posição é a última coluna (linhas 1–3).A orientação do efetuador é a submatriz 3×3 superior–esquerda de $\mathbf{T}^{0}_{E}$, e a posição é o vetor na última coluna (linhas 1–3):
onde:
\begin{equation}
{}^{0}\!\mathbf{R}_{E} = {}^0T_{E}(1\!:\!3,\,1\!:\!3),
\qquad
\vec p_{E} = {}^0T_{E}(1\!:\!3,\,4).
\end{equation}

em que ${}^0T_{E}$ representa a matriz homogênea que fornece a posição e orientação do efetuador final em relação à base da espaçonave.

\subsection{Tabela DH}
A tabela de parâmetros DH considera a geometria do manipulador robótico de configuração 4R1P.
A tabela \eqref{tab:DH-MR} permite calcular as matrizes de transformação e para realizar a análise do espaço de trabalho.

% \begin{table}[ht]
% \centering
% \caption{Parâmetros clássicos de DH}
% \begin{tabular}{c c c c r r}
% \hline
% $i$ & Tipo & $\theta_i$ & $d_i$ [cm] & $a_i$ [cm] & $\alpha_i$ [°] \\
% \hline
% 1 & R & $\theta_1$ & 10                 & 0.0 & 90 \\
% 2 & R & $\theta_2$ & 30                 & 5.5 & 0  \\
% 3 & R & $\theta_3$ & 30                 & 5.5 & 0  \\
% 4 & R & $\theta_4$ & 30                 & 5.5 & 0  \\
% 5 & P & 0.0        & $d_5 [0,15]$   & 0.0 & 0 \\
% \hline
% \end{tabular}
% \label{tab:DH-MR}
% \end{table}

\begin{table}[ht]
\centering
\caption{Parâmetros de Denavit--Hartenberg do MR 4R1P}
\begin{tabular}{c c c c c c}
\hline
$i$ & Tipo & $\theta_i$ & $d_i$ & $a_i$ & $\alpha_i$ \\
\hline
1 & Revoluta & $\theta_1$ & $d_1$ & $a_1$ & $\alpha_1$ \\
2 & Revoluta & $\theta_2$ & $d_2$ & $a_2$ & $\alpha_2$ \\
3 & Revoluta & $\theta_3$ & $d_3$ & $a_3$ & $\alpha_3$ \\
4 & Revoluta & $\theta_4$ & $d_4$ & $a_4$ & $\alpha_4$ \\
5 & Prismática & $0$      & $d_5$ & $a_5$ & $\alpha_5$ \\
\hline
\end{tabular}
\label{tab:DH-MR}
\end{table}
Os tipos de juntas adotadas neste trabalho foram as juntas revolutas (R) e prismática (P), sendo que na junta prismática, a variável é $d_5$; ${}^0 \theta_5$, $a_5$ e $\alpha_5$ são constantes geométricas.
% ---------------------------------------------------------
% Comentário:
% - $d_1 = 10$ cm corresponde à altura entre o centro de massa da base e a junta 1.
% - Os valores $d_2, d_3, d_4 = 30$ cm correspondem aos comprimentos dos elos.
% - Os offsets $a_i = 5.5$ cm representam o deslocamento lateral entre os eixos sucessivos.
% - $\alpha_1 = 90°$ garante a rotação correta do sistema de coordenadas da junta 1 em relação à base.
% - $q_5$ representa o curso prismático (0 a 15 cm) da ferramenta.
% ---------------------------------------------------------

\subsection{Cinemática Direta}

A partir da Tabela DH, a cinemática direta pode ser obtida pela multiplicação sucessiva das matrizes homogêneas
\eqref{eq:matrizes_homogeneas_sucessivas},
onde ${}^0\,\mathbf{T}_{E} \in \mathbb{R}^{4 \times 4}$,
e contém a posição do efetuador final $\{E\}$ em relação à base e a orientação do efetuador.
Sendo assim, a posição do efetuador $\{E\}$ pode ser escrita como:
\begin{equation}
\vec{p}_E = 
\begin{bmatrix}
x_E \\ y_E \\ z_E
\end{bmatrix}.
\end{equation}

Reescrevendo-a, obtém-se:
\begin{equation}
\mathbf{T}^{0}_{E} =
\begin{bmatrix}
{}^{0}\!\mathbf{R}_{E} & \vec{p}_E \\
\mathbf{0}_{1\times 3} & 1
\end{bmatrix}.
\end{equation}

onde ${}^0 \mathbf{R}_{E} \in SO(3)$ representa a orientação do efetuador e $\vec{p}_E \in \mathbb{R}^3$ sua posição.
% ---------------------------------------------------------
% Comentário:
% - A cinemática direta responde à pergunta: dada uma configuração articular, qual é a posição e orientação do efetuador final?
% - É a base para análise de workspace, trajetória e controle.
% ---------------------------------------------------------

\subsection{Cinemática Inversa}
\label{sec:cinematica_inversa}

A cinemática inversa consiste em determinar as variáveis articulares 
$\{ \theta_1, \theta_2, \theta_3, \theta_4, d_5 \}$ 
a partir de uma pose desejada do efetuador final, representada por ${}^0T_E \in SE(3)$. 

A formulação segue de:
\begin{equation}
{}^0 T_E = A_1(\theta_1)\,A_2(\theta_2)\,A_3(\theta_3)\,A_4(\theta_4)\,A_5(d_5),
\label{eq:cinematica_inversa_def}
\end{equation}
onde $A_i$ são as transformações homogêneas DH.
A resolução consiste em inverter \eqref{eq:cinematica_inversa_def}, obtendo:

\begin{equation}
\vec p_E = f_p(\theta_1,\theta_2,\theta_3,\theta_4,d_5), 
\qquad
{}^0R_E = f_R(\theta_1,\theta_2,\theta_3,\theta_4).
\end{equation}

As equações não possuem solução única, podendo admitir múltiplas ou nenhuma configuração articular. 
A resolução pode ser conduzida por métodos analíticos quando viável, ou por métodos numéricos iterativos em casos gerais.

% --------------------------------------------------------
% Comentário:
% - Em manipuladores espaciais, a cinemática inversa é mais complexa:
%   * Pode ter múltiplas soluções.
%   * Pode não ter solução (posição fora do alcance).
%   * Requer escolha da solução que evite colisões ou singularidades.
% - Em aplicações espaciais, soluções numéricas (iterativas) são comuns,
%   devido à geometria não trivial e às restrições de segurança.
% ---------------------------------------------------------
\subsection{Espaço de Trabalho}

O espaço de trabalho (\emph{workspace}) é o volume de pontos no espaço que a ponta da ferramenta $\{E\}$ consegue alcançar quando as juntas se movem dentro de seus limites mecânicos.
Sendo assim, um dado conjunto de limites de juntas e um modelo de cinemática direta, o \emph{workspace} resulta do conjunto de posições possíveis da ferramenta quando o manipulador varre esses limites.
% O espaço de trabalho pode ser classificado como:
% \begin{itemize}
%     \item \textbf{Espaço alcançável}:
%     corresponde ao conjunto de pontos que a ferramenta consegue atingir em pelo menos uma orientação possível;
%     \item \textbf{Espaço de destreza}:
%     subconjunto do espaço de alcance, no qual, a ferramenta pode assumir todas as orientações possíveis em torno de um ponto.
% \end{itemize}

O \emph{workspace} depende da geometria do MR, como:
\begin{enumerate}[label=\roman*)]
    \item comprimentos dos elos ($a_i$);
    \item limites de rotação das juntas revolutas e de deslocamento da junta prismática ($d_5$);
    \item intervalos permitidos para as variáveis articulares ($\theta_1,\theta_2,\theta_3,\theta_4,d_5$);
    \item além das restrições de colisão com a própria estrutura e com o entorno.
\end{enumerate}

% Neste trabalho assume-se base fixa para referencial da geometria do MR, sendo que o efeito base--manipulador (base livre) é tratado adiante.
% ---------------------------------------------------------
% Comentário:
% - A análise do workspace mostra:
%   * A região de captura possível do satélite alvo.
%   * As limitações impostas pelos comprimentos dos elos e cursos das juntas.
%   * Se o manipulador atende ao requisito de aproximação em microgravidade.
% - No contexto de operações espaciais, o workspace deve ser validado
%   em relação à geometria da espaçonave alvo e à zona de exclusão (KOZ).
% ---------------------------------------------------------